% =======================================================
% 4. REQUISITI
% =======================================================
\section{Requisiti}
I requisiti individuati sono classificati e identificati da un codice univoco nel formato \textbf{R[ID]-[Tipo]-[Importanza]}:
\begin{itemize}
	\item \textbf{ID}: Identificativo numerico progressivo.
	\item \textbf{Tipo}: \textbf{F} (Funzionale), \textbf{Q} (Qualità), \textbf{V} (Vincolo), \textbf{P} (Prestazionale).
	\item \textbf{Importanza}: \textbf{O} (Obbligatorio), \textbf{D} (Desiderabile), \textbf{Op} (Opzionale).
\end{itemize}

\subsection{Requisiti funzionali}
% TABELLA LUNGA PER REQUISITI CHE SUPERA LA PAGINA
\renewcommand{\arraystretch}{1.5}
\begin{longtable}{|p{2.5cm}|p{8.5cm}|p{3.5cm}|}
	\caption{Requisiti Funzionali} \\
	\hline
	\rowcolor{primarycolor!10}
	\textbf{Codice} & \textbf{Descrizione} & \textbf{Fonti} \\
	\hline
	\endfirsthead
	
	\hline
	\rowcolor{primarycolor!10}
	\textbf{Codice} & \textbf{Descrizione} & \textbf{Fonti} \\
	\hline
	\endhead
	
	% INSERIRE QUI I REQUISITI
	\textbf{R1-F-O} & [-] & [-] \\ \hline
	\textbf{R2-F-O} & [-] & [-] \\ \hline
	\textbf{R3-F-D} & [-] & [-] \\ \hline
	
\end{longtable}

\subsection{Requisiti di qualità}
\begin{longtable}{|p{2.5cm}|p{8.5cm}|p{3.5cm}|}
	\caption{Requisiti di Qualità} \\
	\hline
	\rowcolor{primarycolor!10}
	\textbf{Codice} & \textbf{Descrizione} & \textbf{Fonti} \\
	\hline
	\endfirsthead
	
	\hline
	\rowcolor{primarycolor!10}
	\textbf{Codice} & \textbf{Descrizione} & \textbf{Fonti} \\
	\hline
	\endhead
	
	\textbf{R1-Q-O} & [-] & [-] \\ \hline
	\textbf{R2-Q-O} & [-] & [-] \\ \hline
\end{longtable}

\subsection{Requisiti di vincolo}
\begin{longtable}{|p{2.5cm}|p{8.5cm}|p{3.5cm}|}
	\caption{Requisiti di Vincolo} \\
	\hline
	\rowcolor{primarycolor!10}
	\textbf{Codice} & \textbf{Descrizione} & \textbf{Fonti} \\
	\hline
	\endfirsthead
	
	\hline
	\rowcolor{primarycolor!10}
	\textbf{Codice} & \textbf{Descrizione} & \textbf{Fonti} \\
	\hline
	\endhead
	
	\textbf{R1-V-O} & [-] & [-] \\ \hline
	\textbf{R2-V-O} & [-] & [-] \\ \hline
\end{longtable}

\subsection{Requisiti prestazionali}
\begin{longtable}{|p{2.5cm}|p{8.5cm}|p{3.5cm}|}
	\caption{Requisiti Prestazionali} \\
	\hline
	\rowcolor{primarycolor!10}
	\textbf{Codice} & \textbf{Descrizione} & \textbf{Fonti} \\
	\hline
	\endfirsthead
	
	\hline
	\rowcolor{primarycolor!10}
	\textbf{Codice} & \textbf{Descrizione} & \textbf{Fonti} \\
	\hline
	\endhead
	
	\textbf{R1-P-O} & [-] & [-] \\ \hline
\end{longtable}